\section{ماژول tt\_compressor (Tensor Train Compressor)}
\label{sec:module_tt_compressor}

\subsection{معرفی و انگیزه}
این مکانیزم بر مبنای مبانی نظری پیشرفته در ادبیات توجه و فشرده‌سازی \cite{ref_ttcompressor_2024} توسعه یافته است.

تجزیه قطار تانسوری (\lr{Tensor-Train Decomposition - TT}) یکی از تکنیک‌های پیشرفته در جبر تانسوری چندبعدی است که برای فشرده‌سازی تانسورهای مرتبه بالا و کاهش پیچیدگی نمایی ابعاد به کار می‌رود. ایده بنیادی این روش، نگاشت یک ماتریس دوبعدی بزرگ به فضایی چندبعدی و سپس نمایش آن به شکل ضرب زنجیره‌ای از هسته‌های تانسوری کم‌مرتبه سه‌بعدی و چهاربعدی است.

موتور فشرده‌سازی \lr{TT Compressor} با هدف جایگزینی لایه‌های متراکم خطی \lr{nn.Linear} و \lr{Conv1D} با ماژول‌های عامل‌بندی‌شده تانسوری به نام \lr{TTLinearLayer} طراحی گردید. با این وجود، تحلیل‌های عمیق تئوری و تجربی در این پروژه نشان داد که اعمال تجزیه قطار تانسوری بر روی ماتریس‌های تبدیل خطی دوبعدی در مدل‌های زبانی ترنسفورمر دارای چالش‌های بنیادین است و به دلیل پدیده درهم‌آمیختگی ابعاد و افت کند مقادیر مقیاس تانسوری نسبت به \lr{SVD} دوبعدی، این ماژول به عنوان یک رویکرد منسوخ‌شده (\lr{Deprecated}) رده‌بندی شده است.

\subsection{مبانی نظری و فرمول‌بندی ریاضی}
فرض کنید ماتریس وزن $\mathbf{W} \in \mathbb{R}^{d_{\text{out}} \times d_{\text{in}}}$ داده شده باشد. ابعاد ورودی و خروجی ابتدا به $d$ عامل صحیح افراز می‌شوند:
\begin{equation}
d_{\text{in}} = \prod_{k=1}^d p_k, \quad d_{\text{out}} = \prod_{k=1}^d q_k
\end{equation}
ماتریس $\mathbf{W}$ به صورت یک تانسور مرتبه $2d$ با ابعاد $(q_1, \dots, q_d, p_1, \dots, p_d)$ بازآرایی و جایگشت می‌یابد تا جفت‌های ورودی-خروجی در کنار یکدیگر قرار گیرند. سپس این تانسور به فرم قطار تانسوری تجزیه می‌شود:
\begin{equation}
\mathbf{W}(i_1, \dots, i_d, j_1, \dots, j_d) = \mathcal{G}_1[i_1, j_1] \cdot \mathcal{G}_2[i_2, j_2] \cdots \mathcal{G}_d[i_d, j_d]
\end{equation}
که در آن $\mathcal{G}_k \in \mathbb{R}^{r_{k-1} \times q_k \times p_k \times r_k}$ هسته‌های تانسوری هستند و رتبه‌های مرزی برابر با $r_0 = r_d = 1$ می‌باشند. مقادیر $\{r_k\}$ به عنوان رتبه‌های قطار تانسوری (\lr{TT-ranks}) شناخته می‌شوند.

تعداد کل پارامترهای ذخیره‌شده در هسته‌های تانسوری برابر است با:
\begin{equation}
N_{\text{TT}} = \sum_{k=1}^d r_{k-1} \cdot p_k \cdot q_k \cdot r_k
\end{equation}
برای دستیابی به کاهش واقعی حجم پارامترها نسبت به ماتریس اولیه ($d_{\text{in}} \cdot d_{\text{out}}$)، رتبه تانسوری باید اکیداً از حد بحرانی کم‌تر باشد:
\begin{equation}
r_{\text{max}} < \frac{d_{\text{out}} \cdot d_{\text{in}}}{\sum_{k=1}^d p_k \cdot q_k}
\end{equation}
در حین گذر پیش‌رو، حاصل‌ضرب تانسور ورودی $\mathbf{X}$ با هسته‌های تانسوری با انقباض متوالی تانسورها (\lr{Tensor Contraction}) با استفاده از دستورات \lr{torch.einsum} یا تبدیل ضرب‌های ماتریسی پی‌درپی محاسبه می‌گردد.

\paragraph{تحلیل طیفی و نامساوی‌های پایداری خطا}
برای بررسی میزان دقت تقریب ماتریس وزن یا نگاشت توجه در ماژول \lr{tt\_compressor}، تفاضل ماتریس فاکتورگیری‌شده از نگاشت متراکم اصلی بر اساس نامساوی هافدینگ-فون‌نیومان ارزیابی می‌گردد. فرض کنید $\mathbf{W} \in \mathbb{R}^{d_{\text{in}} \times d_{\text{out}}}$ ماتریس وزن هدف باشد که توسط تقریب کم‌مرتبه یا پراکنده با رتبه موثر $r$ جایگزین شده است. کران بالای خطای بازسازی بر مبنای نرم فروبنیوس به صورت زیر تصریح می‌شود:
\begin{equation}
\label{eq:spectral_bound_tt_compressor}
\| \mathbf{W} - \mathbf{W}_{r} \|_F = \sqrt{\sum_{i=r+1}^{\min(d_{\text{in}}, d_{\text{out}})} \sigma_i^2(\mathbf{W})} \le \epsilon \| \mathbf{W} \|_F
\end{equation}
که در آن $\sigma_i(\mathbf{W})$ بیانگر $i$-امین مقدار منفرد مرتب‌شده ماتریس به صورت نزولی است. نسبت حفظ انرژی طیفی $\eta_r = \frac{\sum_{i=1}^r \sigma_i^2}{\sum_{i=1}^{d} \sigma_i^2}$ تضمین می‌کند که با انتخاب آستانه $\tau \ge 0.90$، بیش از نود درصد از پراکندگی واریانس داده‌ها حفظ شده و افت سیگنال در انتشار گرادیان رو به عقب به حداقل ممکن کاهش می‌یابد.

\paragraph{تحلیل مجانبی و پیچیدگی اشتراک حافظه سخت‌افزاری}
از منظر پیچیدگی محاسباتی و انتقال داده‌ها در سلسومراتب حافظه (حافظه سراسری \lr{HBM} به حافظه محلی \lr{SRAM})، عملیات ضرب ماتریسی در این مکانیزم نیازمند تخصیص بافر موقت با ابعاد فاکتورگیری‌شده است. برای پردازش دنباله‌ای به طول $N$ با بعد مخفی $D$، بار محاسباتی هسته از فرم درجه دو به رابطه زیر تقلیل می‌یابد:
\begin{equation}
\label{eq:complexity_tt_compressor}
\mathcal{C}_{\text{tt\_compressor}} = \mathcal{O}\left( B \cdot N \cdot D \cdot r + B \cdot H \cdot N \cdot d_k \right) \ll \mathcal{O}\left( B \cdot H \cdot N^2 \cdot d_k \right)
\end{equation}
این بهبود مرتبه پیچیدگی، با کاهش چشمگیر حجم تبادل تانسورها در گذرگاه‌های ارتباطی کارت گرافیک، مانع از اشباع پهنای باند حافظه (پدیده \lr{Memory-Bound}) در فرآیند استنتاج طولانی‌مدت می‌گردد.

\subsection{معماری و پیاده‌سازی}
پیاده‌سازی ماژول در فایل \lr{Codes/framework/compression\_engine/tt\_compressor.py} سازمان‌یافته و شامل کلاس‌های \lr{TTLinearLayer} و \lr{TTCompressor} است که با شناسه \lr{tt} در رجیستری ذخیره شده‌اند:
\begin{itemize}
    \item \textbf{عامل‌بندی خودکار ابعاد (\lr{Dimension Factorization}):} تابع \lr{factorize\_dimension\_pair} ابعاد دلخواه ماتریس را به زوج‌فاکتورهای بهینه با فاصله کمینه از ریشه دوم تجزیه می‌کند تا هسته‌ها متعادل باشند.
    \item \textbf{انطباق با \lr{Conv1D} و شبکه‌های مدرن:} پشتیبانی کامل از ویژگی‌های اختصاصی لایه‌های کانولوشن یک‌بعدی \lr{GPT-2} از طریق نگاشت ویژگی‌های \lr{nf} و \lr{nx}.
    \item \textbf{حفاظت از لایه‌های حساس:} مستثنی کردن لایه‌های پایانی سر دسته‌بندی و واژگان به منظور حفظ پایداری محاسباتی.
    \item \textbf{تعبیه رتبه بیشینه ایمن (\lr{Safe Max-Rank Guard}):} سنجش مرز صرفه‌جویی حافظه با احتساب حاشیه ایمنی \lr{safety\_margin=0.9} برای جلوگیری از افزایش ناخواسته پارامترها در رتبه‌های بالا.
\end{itemize}

\subsection{نتایج تجربی}
نتایج حاصل از ارزیابی‌های تجربی ثبت‌شده برای فشرده‌ساز \lr{tt\_compressor} بر روی مدل‌های بنچ‌مارک در جدول \ref{tab:benchmark_tt_compressor} درج شده است:

\begin{table}[H]
\centering
\small
\caption{نتایج تجربی و اعتبارسنجی بنچ‌مارک فشرده‌ساز \lr{tt\_compressor} روی مدل‌های پایه}
\label{tab:benchmark_tt_compressor}
\begin{tabular}{lccccc}
\toprule
\textbf{مدل پایه} \lr{(Base Model)} & \textbf{دقت} \lr{(Accuracy)} & \textbf{شاخص} \lr{(F1-Score)} & \textbf{حافظه} \lr{(Peak VRAM)} & \textbf{توان} \lr{(Throughput)} & \textbf{وضعیت} \\
\midrule
\lr{DistilBERT} & \lr{59.06\%} & \lr{0.7027} & \lr{435.53 MB} & \lr{69,569.97 tok/s} & \lr{PASS} \\
\lr{RoBERTa} & \lr{50.92\%} & \lr{0.6748} & \lr{642.71 MB} & \lr{7,882.51 tok/s} & \lr{PASS\_LOW} \\
\lr{GPT-2} & \lr{0.00\%} & \lr{0.0000} & \lr{0.00 MB} & \lr{0.00 tok/s} & \lr{CRASH} \\
\lr{Qwen} & \lr{---} & \lr{---} & \lr{---} & \lr{---} & \lr{NOT\_RUN} \\
\bottomrule
\end{tabular}
\end{table}

همان‌گونه که در جدول مشخص است، فشرده‌ساز \lr{tt\_compressor} در مدل \lr{DistilBERT} دقت \lr{59.06\%} و شاخص F1 معادل \lr{0.7027} را به دست آورده است. در مدل \lr{RoBERTa}، به علت افت عملکرد و خطای نسبی بازسازی فروبنیوس بالا، دقت به \lr{50.92\%} افت کرده و وضعیت \lr{PASS\_LOW} ثبت گردیده است. در مدل \lr{GPT-2}، اجرای آزمون به دلیل بروز خطای نرم‌افزاری متغیر درون‌حلقه‌ای (\lr{UnboundLocalError: cannot access local variable 'majority\_count'}) متوقف شده و وضعیت \lr{CRASH} گزارش شده است. این ماژول بر روی مدل \lr{Qwen} ارزیابی نشده است (\lr{NOT\_RUN}).

\subsection{محدودیت‌ها و تحلیل شکست}
دلایل نظری و فنی عدم موفقیت این ماژول و خروج آن از صف بنچ‌مارک فعال به شرح زیر است:
\begin{itemize}
    \item \textbf{درهم‌آمیختگی حالت‌ها و کندی فروپاشی مقادیر مقیاس (\lr{Mode Scrambling}):} ماتریس‌های وزن شبکه‌های عصبی ذاتاً ساختاری دوگانه میان ورودی و خروجی دارند. نگاشت آن‌ها به فضاهای چندبعدی بالاتر موجب برهم‌خوردن پیوستگی داده‌ها و توزیع غیریکنواخت مقادیر منفرد شده و برای حفظ خطای پایین، رتبه‌های \lr{TT} باید بسیار بزرگ انتخاب شوند که مزیت فشرده‌سازی را زایل می‌کند.
    \item \textbf{فقدان صرفه‌جویی واقعی در پارامترها و \lr{FLOPs}:} بر خلاف \lr{SVD} که کاهش ابعاد مستقیم ایجاد می‌کند، انقباض متوالی هسته‌های تانسوری نیازمند عملیات محاسباتی تکراری متعددی است که نه تنها سرعت را بهبود نمی‌بخشد بلکه در ابعاد استاندارد مدل‌های زبانی، کاهش حجم قابل ملاحظه‌ای ایجاد نمی‌کند.
    \item \textbf{خطای ناپایداری در فرآیند ارزیابی کدها:} وابستگی محاسبات پس‌آزمون به متغیرهای محلی محاسبه اکثریت پاسخ در فرآیند استنباط توکن‌های \lr{GPT-2} منجر به شکست اجرایی در اسکریپت‌های ارزیابی شد.
\end{itemize}
